\documentclass[11pt]{article}
\usepackage[T1]{fontenc}
\usepackage{lmodern,microtype,amsmath,amssymb,amsthm,mathtools,mathrsfs,tikz-cd}
\usepackage[margin=1in]{geometry}
\usepackage{xurl,hyperref}
\hypersetup{colorlinks=true,linkcolor=blue,citecolor=blue,urlcolor=blue,
pdftitle={Range completion and computational cosheaves on a three-point space},
pdfauthor={Florian Lengyel}}
\newtheorem{theorem}{Theorem}[section]
\newtheorem{proposition}[theorem]{Proposition}
\newtheorem{lemma}[theorem]{Lemma}
\newtheorem{corollary}[theorem]{Corollary}
\newcommand{\N}{\mathbb N}
\newcommand{\Set}{\mathbf{Set}}
\newcommand{\ran}{\operatorname{ran}}
\newcommand{\Par}{\operatorname{Par}}
\newcommand{\Tot}{\operatorname{Tot}}
\newcommand{\E}{\mathscr E}
\newcommand{\A}{\mathscr A}
\newcommand{\C}{\mathscr C}
\title{Range completion and computational cosheaves\\on a three-point space}
\author{Florian Lengyel}
\date{29 October 2026}
\begin{document}
\maketitle
\begin{abstract}
Free range completions exist, but their universal properties must be
distinguished from preservation of a universal evaluator. We state the
relevant distinction and prove an obstruction for completions that add
partial arrows without adding total arrows.
We then construct a recursion category from computational cosheaves on
a three-point space with a nontrivial two-open cover. Global values are
pushouts presented by computably enumerable equivalence relations.
Finite limits, ranges, and word objects are constructed on the three
basic opens. A countable family of complete bipartite components supplies
a universal object, Heller frames, and a filtered Turing evaluator.
The finite-proof construction therefore gives essential incompleteness
in this category without decidable atom equality.
\end{abstract}

\section{Adjoining ranges and retaining computation}

Composition is written $gf=g\circ f$. In a restriction category,
$\bar f:X\to X$ denotes the domain idempotent of $f:X\to Y$.
The restriction order is $f\leq g$ if $f=g\bar f$;
we write $\boxtimes$ for binary restriction products.
A range idempotent is written $\widehat f:Y\to Y$.
Restriction functors preserve the domain operation; range functors
preserve both operations.

There is a general free range construction. Cockett, Guo, and Hofstra
construct a functor $R$ left adjoint to the forgetful functor from
range categories to restriction categories
\cite[Theorem~5.12]{CGH1}. For a small restriction category $\C$,
its unit $\eta:\C\to R(\C)$ has the universal property
\[
 \operatorname{RangeFun}(R(\C),\mathscr D)
 \cong \operatorname{RestrFun}(\C,\mathscr D)
\]
for every range category $\mathscr D$, naturally in both categories.
The construction keeps the objects and represents arrows by labeled
trees together with an original arrow \cite[\S5]{CGH1}.

Their basic algebraic range axioms do not include
\begin{equation}\label{eq:agreement}
 vf=wf\quad\Longrightarrow\quad
 v\widehat f=w\widehat f.
\end{equation}
This is their additional axiom RR.5 \cite[\S4.11]{CGH1};
it is the universal-agreement property required in \cite{Iteration}.
Here $f:X\to Y$ and $v,w:Y\to Z$ may all be partial.
Thus the existence of the general left adjoint alone does not discharge
that hypothesis.

For a discrete cartesian restriction category, there is also a faithful
discrete cartesian range completion $R_d(\C)$, universal for
product-preserving functors into discrete cartesian range categories
\cite[Theorem~3.27]{CGH2}. Here discrete means that each diagonal has a
restriction inverse; it does not mean Boolean-decidable equality.
In this cartesian setting discreteness is equivalent to
\eqref{eq:agreement} \cite[Proposition~3.11]{CGH2}, so this completion
does supply universal agreement.
The regularization by fractions in \cite[Theorem~3.19]{CGH2} is a
further construction. Regularity, coproducts, word-length coverage,
and universal evaluation must be checked at the relevant step.

\begin{lemma}\label{lem:no-new-total}
Every arrow $q:A\to B$ of $R_d(\C)$ is a restriction of
$\eta f$ for some $f:A\to B$ of $\C$.
In particular, every total arrow of $R_d(\C)$ is the image of a
total arrow of $\C$.
\end{lemma}
\begin{proof}
The span presentation in \cite[Construction~3.21 and Lemma~3.22]{CGH2}
represents $q$ by $[k,h]$, where $k:Z\to A$, $h:Z\to B$,
and $h=fk$ for some $f:A\to B$; one may arrange
$\bar k=\bar h$.
The extension formula in the proof of their Theorem~3.27, applied
to $\eta$ itself, gives the identity extension. Indeed, the
universal equivalence of functor categories supplies an isomorphism
to the identity extension restricting to the identity on $\eta$.
Since $\eta$ is identity on objects, every component of this
isomorphism is an identity. Thus
\[
 q=(\eta f)\widehat{\eta k}.
\]
Since $\widehat{\eta k}$ is a restriction idempotent, $q\leq\eta f$.
If $q$ is total, this inequality gives $q=\eta f$.
Faithfulness of $\eta$ and preservation of restriction then give
$\bar f=1_A$.
\end{proof}

The following elementary obstruction concerns preservation of the
\emph{given} evaluator. Its assumptions are stated independently of
any particular completion.

\begin{proposition}\label{prop:evaluator-obstruction}
Let $J:\C\to\mathscr D$ be a restriction functor preserving binary
restriction products. Suppose that every total arrow
$JX\to JU$ in $\mathscr D$ is $Jc$ for some total $c:X\to U$ in
$\C$. Suppose also that every object $Y$ of $\C$ has a total
restriction embedding $j_Y:Y\to U$ with partial retraction
$p_Y:U\to Y$.
If $\tau:U\boxtimes U\to U$ is an arrow of $\C$ and
$J\tau$ is parameter-universal for arrows
$JX\boxtimes JU\to JU$ in $\mathscr D$, then $J$ is full on
arrows between its image objects.
\end{proposition}
\begin{proof}
We identify preserved products through the canonical product comparisons.
Let $g:JX\to JY$ be partial. Put
\[
 G=(Jj_Y)g\pi_1:JX\boxtimes JU\to JU.
\]
Universality supplies a total $d:JX\to JU$ with
$G=(J\tau)(d\boxtimes1)$. Write $d=Jc$ with $c$ total.
Then $G=JG_0$, where $G_0=\tau(c\boxtimes1_U)$.
Precompose with $J\langle1_X,j_X\rangle$ and postcompose with
$Jp_Y$. Since $p_Yj_Y=1_Y$ and
$\pi_1\langle1_X,j_X\rangle=1_X$, this gives
\[
 g=J\bigl(p_YG_0\langle1_X,j_X\rangle\bigr).
\]
The equations include their restriction domains.
\end{proof}

Consequently a completion satisfying the displayed condition on total
arrows cannot both add genuinely new partial arrows between old objects
and preserve the original evaluator's universality.
By Lemma~\ref{lem:no-new-total}, this applies to $R_d(\C)$ whenever
the original objects restriction-embed in $U$: if its unit is not full,
$\eta\tau$ fails parameter-universality.
A replacement
evaluator or a completion adding new total program maps is a different
question. This observation does not rule out a computational range
completion; it identifies a property such a construction must address.

\section{Cosheaves and a computational presentation}\label{sec:cosheaves}

Let $X=\{a,b,c\}$ have opens
\[
 \varnothing,\quad D=\{c\},\quad
 L=\{a,c\},\quad R=\{b,c\},\quad X.
\]
Write $\mathcal O(X)$ for the poset category of opens and inclusions.
A set-valued cosheaf is a covariant functor
$C:\mathcal O(X)\to\Set$ with $C(\varnothing)=\varnothing$
whose cover diagram
\[
 \coprod_{i,j}C(V_i\cap V_j)\rightrightarrows
 \coprod_iC(V_i)\longrightarrow C(V)
\]
is a coequalizer whenever the $V_i$ cover $V$
\cite[\S\S2.1--2.2]{Curry}. Arrows are natural transformations.
Let $J_0$ be the category with objects $0,1,2$ and two
nonidentity arrows $0\to1$, $0\to2$.

\begin{proposition}\label{prop:cosheaf-equivalence}
Restriction to $D,L,R$ gives an equivalence
\[
 \operatorname{CoSh}(X,\Set)\simeq\Set^{J_0}.
\]
Its inverse sends a span $A\xleftarrow{r}C\xrightarrow{s}B$
to the cosheaf with values $C,A,B$ on $D,L,R$ and
$A\amalg_C B$ on $X$.
\end{proposition}
\begin{proof}
Every cover of $D,L$, or $R$ contains the covered open itself.
A cover of $X$ either contains $X$ or contains both $L$ and $R$.
Thus its only additional condition is the pushout square
\[
\begin{tikzcd}
C\arrow[r,"s"]\arrow[d,"r"']&B\arrow[d]\\
A\arrow[r]&A\amalg_C B.
\end{tikzcd}
\]
Redundant opens in a cover give the same coequalizer, since their
maps already factor through the two displayed cover members.
A commuting map of spans induces a unique map of pushouts.
This defines the inverse functor on arrows; the pushout universal
property gives both inverse natural isomorphisms.
\end{proof}

Let $\C_{\rm rec}$ have c.e.\ subsets of finite powers of $\N$
as objects and partial recursive functions total on their sources as
arrows; equality of arrows is extensional. Here c.e.\ means
computably enumerable. Define
\[
 \E=\C_{\rm rec}^{J_0},\qquad \A=\Par(\E,\operatorname{Mono}).
\]
Here $\Par(\E,\operatorname{Mono})$ has the objects of $\E$;
an arrow $Y\to Z$ is an isomorphism class of spans
$Y\xleftarrow{m}T\xrightarrow{f}Z$ with $m$ monic.
Composition uses pullback, and total arrows are represented by $m=1$.
The required finite limits are verified below.
An object of $\E$ is a pair of recursive transitions
$r_i:Y_0\to Y_i$, $i=1,2$, on c.e.\ carriers.
An arrow is a triple $f_j:Y_j\to Z_j$ with
$r_i^Zf_0=f_i r_i^Y$.
We call these computational cosheaves, using
Proposition~\ref{prop:cosheaf-equivalence} for their underlying
set-valued cosheaves. The permitted natural transformations are
exactly those computable at the three basic opens.

For a computational span, its global pushout has an effective quotient
presentation. Take the tagged carrier $Y_1+Y_2$ and the equivalence
relation generated by
\[
 (1,r_1 y)\sim(2,r_2 y)\qquad(y\in Y_0).
\]
Enumeration of finite paths in this c.e.\ edge relation enumerates
the equivalence relation. A computational span map acts on tagged
representatives and respects that relation, so induces the global
map. Composition is induced by composition on representatives.
No decision procedure for quotient equality is required.

This qualification is necessary. If $E$ is any c.e.\ equivalence
relation on $\N$, take $Y_0=E$, $Y_1=Y_2=\N$ and the two
projections. The pushout is $\N/E$: reflexivity joins the two
copies of each integer, and the other edges impose precisely $E$.
An effective injective coding of this quotient into a numerical
carrier, computable from representatives, would decide $E$ by
equality of the codes. Such a coding therefore cannot be demanded
for arbitrary c.e.\ $E$.

\begin{proposition}\label{prop:structural}
The category $\E$ is finitely complete, regular, and extensive.
Its monomorphisms and regular images are computed at $0,1,2$.
The partial category $\A$ is split, has restriction zeros,
restriction products and coproducts, weak choice, and universal
product-compatible ranges. Componentwise nonempty words give
Heller's $B^\#$-structure with stable finite-length coverage.
\end{proposition}
\begin{proof}
In $\C_{\rm rec}$ finite limits use tuples and equality of
numerical codes; images are c.e.\ sets. A total surjection admits a
partial recursive section total on its target: enumerate its source
and search for a preimage. Hence it coequalizes its kernel pair.
Surjections are stable under pullback. Tagged finite sums give
extensivity.

Finite limits, finite coproducts, and image factorizations of
diagrams of shape $J_0$ are computed at the three objects.
The transitions restrict to image carriers by naturality.
For componentwise regular epimorphisms, componentwise descent is
natural: its commuting equations follow after precomposition with
the corresponding epic component. Thus these are regular
epimorphisms in $\E$, and are pullback-stable.
Decomposing each component over a tagged sum gives subdiagrams,
so the slice equivalence for extensivity also holds.

Monicity is detected by the three representable finite diagrams:
$(1\leftarrow1\to1)$, $(1\leftarrow0\to0)$, and
$(0\leftarrow0\to1)$, where $0$ and $1$ denote empty and singleton carriers.
A chosen component element defines a computable map from the
corresponding diagram. Hence a monomorphism is injective at
each basic component. Conversely, these injections imply
monicity. Their inverses on their c.e.\ images are partial recursive.

Consequently a partial arrow in $\A$ is represented by c.e.\
subsets $T_j\subseteq Y_j$ with $r_i(T_0)\subseteq T_i$,
and recursive value maps on the $T_j$ commuting with the
restricted transitions. Composition pulls back the domain subsets.
Restrictions are the partial identities on these subdiagrams,
which split. Products and coproducts use the diagram operations.
Ranges are the three c.e.\ value images. Equality of partial
composites includes equality of domains; componentwise surjectivity
onto the image reflects that equality and equality of values.
This proves universal range agreement. Products preserve these
images.

A partial monomorphism has an injective value map on its domain:
test the defining monicity condition on total maps into that domain.
Reversing its span gives a section, proving weak choice.

Finally, apply nonempty-word formation to each carrier, and apply
each transition to every letter. Concatenation, singleton insertion,
flattening, left and right decomposition, parameter repetition,
and alternating-word decomposition are recursive and natural.
They obey Heller's seven-map axioms componentwise.
The fixed-length inclusions cover all words. Equality on every
length implies equality of partial domains and values, including
after product with any diagram. This is stable agreement coverage,
without an assumption of arbitrary external c.e.\ unions.
\end{proof}

\section{Internal products and domains}\label{sec:internal}

The equivalence in Proposition~\ref{prop:cosheaf-equivalence}
does not make every open-set evaluation preserve finite limits.
For spans $Y,Z$, their internal product has basic components
$Y_j\times Z_j$ and global value
\[
 (Y_1\times Z_1)
 \amalg_{Y_0\times Z_0}
 (Y_2\times Z_2).
\]
Its projections induce the product universal property from that
of diagrams. It need not be the product of the two global
pushouts. For $H=(1\leftarrow0\to1)$, one has $H\times H=H$
but $H(X)\times H(X)$ has four elements, whereas
$(H\times H)(X)$ has two.

Likewise, the componentwise inclusion
\[
 (1\leftarrow0\to1)\ \longrightarrow\ (1\leftarrow1\to1)
\]
is a categorical monomorphism, while its global map is $2\to1$.
Thus a partial map in $\A$ is not generally a family of
pointwise partial functions on all five opens. Its domain is a
categorical subobject of the computational cosheaf, represented
by an injective subdiagram at the three basic opens.
The images in Proposition~\ref{prop:structural} are categorical
regular images, not an assertion about ordinary subsets of global
pushouts.

The same distinction applies to word objects and syntax.
Their basic components are words or formulas in the indicated
carriers; their global values are obtained by the pushout
construction. There is no assertion that this pushout equals
the words or formulas over the global atom quotient.
These internal objects satisfy the categorical universal properties
used below.

\section{A universal object with frames}\label{sec:universal}

Define $V$ by
\[
 V_0=\N^4,\quad V_1=V_2=\N^2,\qquad
 r_1(k,z,a,b)=(k,a),\quad r_2(k,z,a,b)=(k,b).
\]
For each fixed $k$, this is a complete bipartite incidence diagram
with infinitely many central labels over every endpoint pair.
The index $k$ distinguishes its connected components.

\begin{lemma}\label{lem:dominance}
Every object of $\E$ has a total restriction embedding into $V$.
\end{lemma}
\begin{proof}
Choose recursive injections $j_i:Y_i\to\N$ with partial recursive
inverses on their images. Put
\[
 m_i(y_i)=(0,j_i y_i)\quad(i=1,2),\qquad
 m_0(y)=(0,j_0y,j_1r_1y,j_2r_2y).
\]
These are commuting componentwise injections. Their inverse
maps on the image subdiagram give the restriction inverse in $\A$.
\end{proof}

Fix a recursive pairing bijection $\N^2\cong\N$, a recursive
bijection $\N\cong\N^\#$ onto nonempty numerical words, and
recursive bijections $\N\cong\N^n$ uniform for $n\ge1$.

\begin{lemma}\label{lem:frames}
There are total computable isomorphisms
$V\cong V\boxtimes V$, $V\cong V+V$, and $V\cong V^\#$.
\end{lemma}
\begin{proof}
For the product frame, decode $k$ into $(k',k'')$ and independently
decode each of $z,a,b$ into a pair. At endpoint $i$, decode
$(k,x)$ into $((k',x'),(k'',x''))$. The two transition equations
hold coordinatewise; pairing supplies the inverse.

For the coproduct frame, write $k=2q+t$, $t\in\{0,1\}$.
Choose summand $t$, replace $k$ by $q$, and keep the other
coordinates. This is natural and has the evident inverse.

For the word frame, decode $k$ as a nonempty component word
$(k_1,\ldots,k_n)$. Decode each of $z,a,b$ into $n$ coordinates,
giving the word
$((k_j,z_j,a_j,b_j))_{j=1}^n$.
At an endpoint decode $(k,x)$ into $((k_j,x_j))_{j=1}^n$.
Transitions commute with these maps. Recoding the component
word and coordinate tuples gives total recursive inverses.
\end{proof}

The component index is essential for this choice of frames.
The single-component diagram
$\N\leftarrow\N^3\to\N$ is connected: every left endpoint
is joined to every right endpoint. Its coproduct with itself
has two components, so the two diagrams are not isomorphic.
Passing to $V$ supplies enough component indices for all three
frame constructions.

\section{A compatible universal evaluator}\label{sec:evaluator}

Fix a universal partial recursive function $\varphi_p(x)$ with
total recursive specialization. Numerical tuple encodings are
understood below; every code decodes to a tuple of the stated
fixed length.

\begin{theorem}\label{thm:turing}
There is a parameter-universal Turing morphism
$\tau:V\boxtimes V\to V$ in $\A$.
\end{theorem}
\begin{proof}
At endpoint $i=1,2$ define
\[
 \tau_i((k_p,p_i),(k_x,x_i))
   =\operatorname{decode}_2
      \varphi_{p_i}(\operatorname{encode}_2(k_x,x_i)).
\]
The program-component index $k_p$ is ignored.
At component $0$, a program and an input have the forms
$(k_p,p_0,p_1,p_2)$ and $(k_x,x_0,x_1,x_2)$.
Run the central program
\[
 \operatorname{decode}_4
 \varphi_{p_0}(\operatorname{encode}_4(k_x,x_0,x_1,x_2))
       =(\ell,y_0,y_1,y_2).
\]
Return this tuple precisely when both endpoint computations also
halt and return $(\ell,y_1)$ and $(\ell,y_2)$, respectively.
Dovetailing and numerical equality define a partial recursive map.
Its domain projects into both endpoint domains, and its values
satisfy both transition equations. Hence $\tau$ is an arrow of $\A$.

Let $F:Y\boxtimes V\to V$ be partial. Represent its three
components by programs including their c.e.\ domain tests.
Specialization gives total-on-source recursive index maps $s_i$
such that the program $s_i(y_i)$ computes
$x\mapsto F_i(y_i,x)$ with exactly its specified domain.
Define
\[
 c_i(y_i)=(0,s_i(y_i))\quad(i=1,2),\qquad
 c_0(y)=(0,s_0(y),s_1(r_1y),s_2(r_2y)).
\]
Then $c:Y\to V$ is a total commuting diagram.
The equality $\tau_i(c_i,-)=F_i$ holds at both endpoints.
Where the central program for $F_0$ halts, the domain and
naturality conditions for $F$ guarantee that both endpoint
checks succeed. Conversely, the central evaluator cannot halt
without that program halting. Thus
$\tau(c\boxtimes1_V)=F$, with equality of domains at all
three components.
\end{proof}

Let $\mathscr R_X$ be the full partial category on formal objects
generated by $V$ under $\boxtimes,+,\#$.
Lemma~\ref{lem:frames} makes it isotypical.
Theorem~\ref{thm:turing} restricts to these objects, so
$\mathscr R_X$ is a recursion category.
By Lemma~\ref{lem:dominance}, its restriction-idempotent splitting
is equivalent to $\A$. In particular, the ranges in this
example are verified within the computational category.

\section{Equality, case distinction, and finite proofs}

\begin{proposition}\label{prop:distinction}
The object $V$ has no decidable equality in $\E$.
The terminal object of $\E$ has five subobjects up to isomorphism,
so $\E$ is equivalent to neither the numerical total category
nor the computational Sierpi\'nski total category.
\end{proposition}
\begin{proof}
The Boolean object $2=1+1$ is the constant diagram $2$.
If $\epsilon:V\times V\to2$ classified the diagonal, naturality
along $r_1$ would make $\epsilon_0(u,v)$ true whenever
$r_1u=r_1v$. Choose distinct central tuples differing only in
$z$; their endpoint values agree. This contradicts the required
central diagonal inverse image.

A subobject of the terminal diagram has either empty central
component and arbitrary empty/singleton endpoint components
(four choices), or all three components singleton (one choice).
The numerical and Sierpi\'nski terminal subobject orders have
two and three elements, respectively \cite[\S6]{Iteration}.
Equivalences preserve these orders.
\end{proof}

\begin{lemma}\label{lem:case}
There are disjoint domains $q_0,q_1$ on $V$ with the binary
case-distinction property: for every complementary pair of
domains $d,d'$ on $V$, a total $h:V\to V$ satisfies
$h^*q_0=d'$ and $h^*q_1=d$.
\end{lemma}
\begin{proof}
Let $q_i$ be the range of the global point with central value
$(i,0,0,0)$ and endpoint values $(i,0)$, for $i=0,1$.

Write the component subsets of $d,d'$ as complementary partitions.
For every tuple $(k,z,a,b)$, its two endpoints must lie in the
same part: its central part projects into that part at both
endpoints. Varying $a,b$ shows that every endpoint with the
same $k$ has the same part. Every central tuple with that $k$
is forced into that part as well. Thus $d$ is the union of
whole component blocks indexed by a subset $K\subseteq\N$.
The set $K$ and its complement are c.e., by testing the two
domains at the endpoint $(k,0)$. Parallel search decides $K$.
Let $h$ send block $k$ to the displayed point $1$ if $k\in K$
and to point $0$ otherwise. Its three component maps are total
recursive and commute. Their inverse-image domains are as claimed.
\end{proof}

\begin{theorem}\label{thm:essential}
In the restriction splitting of $\mathscr R_X$, there is a
consistent computational classical propositional presentation
such that every consistent computational extension is incomplete.
Its syntax and proof algebras are parameter-stable finite-stage
algebras, and its theorem domain is the selected range of its
proof-conclusion map.
\end{theorem}
\begin{proof}
Proposition~\ref{prop:structural}, Lemma~\ref{lem:dominance},
and Theorem~\ref{thm:turing} verify every hypothesis of
\cite[Corollary~4.2]{Iteration}. The finite trace/worklist
construction therefore supplies both stage unions, and the
typed Turing-fold theorem cited there supplies the folds.
The typing filter is retained before applying the polynomial;
its strength uses only total program-parameter maps.

The object $V$ is nonzero. Lemma~\ref{lem:case} and the
Turing diagonal give disjoint inseparable domains $a,b$ on $V$.
Use syntax $V+Z+Z^2$, the three Hilbert schemes and modus ponens,
and the signed atomic axioms indexed by the split objects of
$a,b$. The equality-constraint argument of
\cite[Theorem~5.3]{Iteration} proves consistency without a
Boolean equality test. Regular images give the least rule-closed
theorem domain. The one-node proof maps give its two atomic
range inequalities, so a complete consistent computational
extension would separate $a,b$ by complementary domains,
contradicting inseparability.
\end{proof}

This verifies a nontrivial-cover cosheaf instance. It does not
assert a general recursion-preserving range completion or closure
of arbitrary Heller-generated subcategories under ranges.
The next completion question is to identify which added range
predicates admit a uniform evaluator while retaining the required
total program maps, word-length coverage, and finite coproducts.
The three-point example supplies a test category in which these
properties can be compared without assuming that global pushouts
have decidable equality.

\begin{thebibliography}{9}
\bibitem{CGH1} J.~R.~B.~Cockett, X.~Guo, and P.~Hofstra,
Range categories I: General theory,
\emph{Theory and Applications of Categories} \textbf{26}(17)
(2012), 412--452.
\url{https://www.tac.mta.ca/tac/volumes/26/17/26-17abs.html}.
\bibitem{CGH2} J.~R.~B.~Cockett, X.~Guo, and P.~Hofstra,
Range categories II: Towards regularity,
\emph{Theory and Applications of Categories} \textbf{26}(18)
(2012), 453--500.
\url{https://www.tac.mta.ca/tac/volumes/26/18/26-18abs.html}.
\bibitem{Iteration} F.~Lengyel,
\emph{Iteration domains, finite proofs, and computational cosheaves},
29 October 2026, repository version \texttt{8961637}.
\url{https://github.com/flengyel/CategoricalEssentialIncompleteness/blob/896163707ed3eaa61a30f1d881e8a9f18a3ec0c8/manuscript/iteration_domains_and_cosheaves.tex}.
\bibitem{Curry} J.~M.~Curry,
\emph{Sheaves, Cosheaves and Applications}, Ph.D. thesis,
University of Pennsylvania, 2014.
\href{https://arxiv.org/abs/1303.3255}{arXiv:1303.3255v2}.
\end{thebibliography}
\end{document}
